% =========================================================
% Rational Cauchy Sequences
% =========================================================

\subsection{Rational Cauchy Sequences}

% ---------------------------------------------------------
% Teacher exposition
% ---------------------------------------------------------

\begin{remark}[Teacher exposition: Why we study Cauchy sequences]
\label{rem:cauchy-teacher}

Up to this point we have constructed the rational numbers
\[
\mathbb{Q}
\]
from the integers.  The rational numbers form a very rich number system:
we can add, subtract, multiply, and divide.

However, an important phenomenon appears when we study sequences of
rational numbers.  Some sequences of rational numbers appear to
approach a limit that is \emph{not rational}.

For example, the sequence

\[
1,\;1.4,\;1.41,\;1.414,\;1.4142,\dots
\]

appears to approach the number $\sqrt{2}$, but $\sqrt{2}$ is not a
rational number.

Thus the rational numbers are missing some limits.  In order to
understand this problem precisely, we introduce the concept of a
\emph{Cauchy sequence}.  The key idea is to detect convergence by
examining how the terms of the sequence behave relative to one another,
rather than referring to a limit directly.

\end{remark}

% ---------------------------------------------------------
% Definition
% ---------------------------------------------------------

\begin{definition}[Sequence of Rational Numbers]
\label{def:rational-sequence}

A \emph{sequence of rational numbers} is a function

\[
x : \mathbb{N} \rightarrow \mathbb{Q}.
\]

It is usually written

\[
(x_n)_{n\in\mathbb{N}}
\]

or simply

\[
x_1,x_2,x_3,\dots
\]

where each $x_n \in \mathbb{Q}$.

\end{definition}

% ---------------------------------------------------------

\begin{remark}[Teacher exposition: What a sequence really is]
\label{rem:sequence-teacher}

Although sequences are often written as lists, it is important to
remember that a sequence is formally a function

\[
\mathbb{N} \to \mathbb{Q}.
\]

The natural number $n$ indicates the position in the sequence, and the
value $x_n$ is the rational number appearing in that position.

This functional viewpoint becomes important when we later manipulate
sequences algebraically.

\end{remark}

% ---------------------------------------------------------
% Definition
% ---------------------------------------------------------

\begin{definition}[Cauchy Sequence in $\mathbb{Q}$]
\label{def:cauchy-rational}

A sequence $(x_n)$ of rational numbers is called a
\emph{Cauchy sequence} if

\[
\forall \varepsilon > 0 \;
\exists N \in \mathbb{N}
\quad
\forall m,n \ge N
\]

\[
|x_m - x_n| < \varepsilon.
\]

\end{definition}

% ---------------------------------------------------------
% Intuition
% ---------------------------------------------------------

\begin{remark}[Intuition]
\label{rem:cauchy-intuition}

A Cauchy sequence is one whose terms become arbitrarily close to one
another.

In other words, the sequence eventually becomes tightly clustered
together.

The definition does not mention a limit; it only describes the
internal behavior of the sequence.

\end{remark}

% ---------------------------------------------------------

\begin{remark}[Teacher exposition: Why the definition works]
\label{rem:cauchy-teacher-2}

Suppose a sequence converges to a number $L$.  Then eventually all
terms of the sequence must lie very close to $L$.  Consequently, two
terms taken far enough out in the sequence must also lie very close to
each other.

The Cauchy condition captures exactly this idea: it says that the
sequence becomes internally consistent, with later terms differing
from each other by arbitrarily small amounts.

The remarkable feature of this definition is that it does not require
knowledge of the limit itself.

\end{remark}

% ---------------------------------------------------------
% Examples
% ---------------------------------------------------------

\begin{remark}[Examples]

\textbf{Example 1.}

\[
x_n = \frac{1}{n}
\]

is a Cauchy sequence in $\mathbb{Q}$ because its terms approach $0$.

\medskip

\textbf{Example 2.}

The decimal approximations of $\sqrt{2}$ form a Cauchy sequence:

\[
1,\;
1.4,\;
1.41,\;
1.414,\;
1.4142,\;\dots
\]

Although $\sqrt{2}$ is not rational, the sequence of rational
approximations becomes arbitrarily close together.

\end{remark}

% ---------------------------------------------------------
% Key property
% ---------------------------------------------------------

\begin{theorem}
\label{thm:convergent-implies-cauchy-q}

Every convergent sequence in $\mathbb{Q}$ is a Cauchy sequence.

\end{theorem}

% ---------------------------------------------------------

\begin{remark}[Teacher exposition: One-way implication]
\label{rem:cauchy-one-way}

If a sequence converges, it must be Cauchy.  Intuitively, once the
terms of a sequence are close to the limit, they must also be close to
each other.

However, the converse statement is not always true in $\mathbb{Q}$.
This asymmetry is the key reason that the rational numbers are not a
complete number system.

\end{remark}

% ---------------------------------------------------------

\begin{remark}[Important Observation]
\label{rem:cauchy-not-convergent}

The converse is not always true in $\mathbb{Q}$.

There exist Cauchy sequences of rational numbers that do not converge
to a rational number.

For example, the rational approximations to $\sqrt{2}$ form a Cauchy
sequence in $\mathbb{Q}$ but do not converge to any rational number.

This phenomenon reveals that the rational numbers are
\emph{incomplete}.

\end{remark}

% ---------------------------------------------------------
% Consequence
% ---------------------------------------------------------

\begin{remark}[Motivation for the Real Numbers]
\label{rem:real-motivation}

The incompleteness of $\mathbb{Q}$ motivates the construction of the
real numbers.

One standard construction defines the real numbers as equivalence
classes of Cauchy sequences of rational numbers.

Two sequences $(x_n)$ and $(y_n)$ are considered equivalent if

\[
\lim_{n\to\infty}(x_n - y_n) = 0.
\]

Symbolically,

\[
\mathbb{R}
=
\{\text{Cauchy sequences of rationals}\}/\sim .
\]

Thus the real numbers are obtained by formally adjoining the limits
that are missing from $\mathbb{Q}$.

\end{remark}

% ---------------------------------------------------------

\begin{remark}[Big Picture]
\label{rem:completion-big-picture}

The construction of the real numbers can be viewed as a process of
\emph{completion}.

\[
\mathbb{Z}
\quad \longrightarrow \quad
\mathbb{Q}
\quad \longrightarrow \quad
\mathbb{R}
\]

The integers are extended to the rational numbers by introducing
fractions.  The rational numbers are then extended to the real numbers
by introducing limits of Cauchy sequences.

In this sense the real numbers fill in the ``gaps'' that exist in
$\mathbb{Q}$.

\end{remark}